\section{Conclusion}

The first-order self-tuning regulator of Example 7.6 is implemented on the
ESP32 rig, with the encoder scaling fixed by the one-revolution test at
\(\mathrm{CPR}=44\). The study explains why this is the order that can work
there: the mechanical gain \(b_1\) is recovered to \(\SI{0.06}{\percent}\)
even at that resolution. The same derivation carried one order higher
identifies all four parameters in simulation at \(T=\SI{2}{ms}\) with an exact
speed signal. On the rig the second-order model is structurally identifiable
but not practically identifiable. At \(T=\SI{80}{ms}\) the electrical mode
contributes about \(\SI{1}{\percent}\) of the output and is lost in the encoder
quantisation; shortening the sampling time restores the sensitivity but the
resolution collapses as \(1/T\), so about \(700\) counts per revolution would
be needed instead of \(44\). The second-order design therefore remains a
simulation result, and the stable closed loop measured on the rig must not be
mistaken for a successful identification.
